% Les marqueurs nominatifs du rapport, et le seul endroit du dépôt où un nom de
% personne peut figurer. Tout autre fichier les emploie par leur commande, jamais
% en écrivant la valeur.
%
% Un marqueur laissé vide fait échouer un contrôle : c'est la propriété que le
% critère de terminaison du projet exige, et tests/test_marqueurs_nominatifs.py
% la vérifie sur CE fichier, jamais sur le PDF produit.
%
% Renseigner une valeur, c'est écrire entre les accolades de la commande. Ne pas
% renommer les commandes : le contrôle les cherche par leur nom.

% --- l'état du document -------------------------------------------------------
% Deux valeurs, et deux seulement.
%
%   brouillon : le rapport n'est pas remis. Les cinq marqueurs nominatifs
%               ci-dessous sont TOUS vides, et la page de garde porte une mention
%               visible de version de travail.
%   remise    : le rapport est remis. Les cinq marqueurs sont TOUS renseignés, et
%               la mention disparaît.
%
% Un contrôle vérifie l'accord entre cet état et les cinq marqueurs : renseigner
% quatre marqueurs sur cinq est rouge dans les deux états, et c'est l'oubli qui
% arrive. Passer à la remise ne coûte que ce mot — et le contrôle dira alors quel
% marqueur manque encore.
\newcommand{\etatDuDocument}{brouillon}

% --- marqueurs nominatifs : ils portent un nom de personne ---------------------
\newcommand{\marqueurAuteur}{}
\newcommand{\marqueurEncadrantAcademique}{}
\newcommand{\marqueurEncadrantProfessionnel}{}
\newcommand{\marqueurPresidentJury}{}
\newcommand{\marqueurExaminateur}{}

% --- marqueurs de contexte : ils ne portent aucun nom de personne --------------
% Le contrôle ne les exige pas renseignés ; ils sont ici pour que la page de
% garde n'écrive aucune valeur en dur.
\newcommand{\marqueurEtablissementFormation}{}
\newcommand{\marqueurFiliere}{}
\newcommand{\marqueurAnneeUniversitaire}{}
\newcommand{\marqueurOrganismeAccueil}{Hôpital Sidi Saïd, Meknès}
\newcommand{\marqueurDateSoutenance}{}

% --- ce que l'état fait porter au document ------------------------------------
% Le mécanisme est typographique : le document se suffit, et ne dépend d'aucun
% contrôle pour dire ce qu'il est. `\siBrouillon{...}` compose son argument en
% état de brouillon, et rien en état de remise.
%
% La comparaison passe par \ifx sur deux commandes développées : c'est la forme
% que le noyau offre sans paquet supplémentaire, et elle compare les définitions
% caractère pour caractère.
\newcommand{\etatBrouillonAttendu}{brouillon}
\newcommand{\siBrouillon}[1]{%
  \ifx\etatDuDocument\etatBrouillonAttendu#1\fi%
}

% La mention elle-même, écrite une fois et employée par les deux documents.
\newcommand{\mentionVersionDeTravail}{%
  \siBrouillon{\textbf{Version de travail} — document non remis}%
}
